% Figure: System Representation Conversions
% Chapter 4: Mathematical Representations
\begin{figure}[htbp]
\centering
\begin{tikzpicture}[
    repr box/.style={draw=UofSCBlack, fill=#1!15, rectangle,
                     minimum height=2cm, minimum width=3.8cm,
                     sharp corners, font=\small, align=center, line width=1pt},
    conv arrow/.style={->, >=stealth, thick, UofSC70Black},
    conv label/.style={font=\tiny, fill=white, inner sep=1pt}
]
    % Four representation boxes arranged in a square

    % Differential Equation (top left)
    \node[repr box=UofSCAtlantic] (de) at (-3.5, 2.5) {
        \textbf{Differential Equation}\\[4pt]
        $a_n\dfrac{d^n y}{dt^n} + \cdots + a_0 y$\\[2pt]
        $= b_m\dfrac{d^m u}{dt^m} + \cdots + b_0 u$
    };

    % Transfer Function (top right)
    \node[repr box=UofSCHorseshoe] (tf) at (3.5, 2.5) {
        \textbf{Transfer Function}\\[4pt]
        $G(s) = \dfrac{Y(s)}{U(s)}$\\[2pt]
        $= \dfrac{b_m s^m + \cdots + b_0}{a_n s^n + \cdots + a_0}$
    };

    % State Space (bottom left)
    \node[repr box=UofSCGarnet] (ss) at (-3.5, -2.5) {
        \textbf{State-Space}\\[4pt]
        $\dot{\mathbf{x}} = \mathbf{A}\mathbf{x} + \mathbf{B}u$\\[2pt]
        $y = \mathbf{C}\mathbf{x} + Du$
    };

    % Block Diagram (bottom right)
    \node[repr box=UofSCCongaree] (bd) at (3.5, -2.5) {
        \textbf{Block Diagram}\\[4pt]
        \begin{tikzpicture}[scale=0.5]
            \node[draw, fill=white, minimum width=1cm, minimum height=0.6cm, font=\tiny] (g) at (0,0) {$G(s)$};
            \draw[->, thick] (-1.2,0) -- (g.west);
            \draw[->, thick] (g.east) -- (1.2,0);
        \end{tikzpicture}
    };

    % Horizontal conversions (top)
    \draw[conv arrow] (de.east) -- (tf.west);
    \node[conv label, above] at (0, 2.7) {Laplace Transform};

    \draw[conv arrow] (tf.west) -- (de.east);
    \node[conv label, below] at (0, 2.3) {Inverse Laplace};

    % Horizontal conversions (bottom)
    \draw[conv arrow] (ss.east) -- (bd.west);
    \node[conv label, above] at (0, -2.3) {$G(s) = \mathbf{C}(s\mathbf{I}-\mathbf{A})^{-1}\mathbf{B} + D$};

    \draw[conv arrow] (bd.west) -- (ss.east);
    \node[conv label, below] at (0, -2.7) {Realization};

    % Vertical conversions (left)
    \draw[conv arrow] (de.south) -- (ss.north);
    \node[conv label, left, align=right] at (-3.7, 0) {State\\selection};

    \draw[conv arrow] (ss.north) -- (de.south);
    \node[conv label, right, align=left] at (-3.3, 0) {Eliminate\\states};

    % Vertical conversions (right)
    \draw[conv arrow] (tf.south) -- (bd.north);
    \node[conv label, right, align=left] at (3.7, 0) {Direct\\mapping};

    \draw[conv arrow] (bd.north) -- (tf.south);
    \node[conv label, left, align=right] at (3.3, 0) {Block\\reduction};

    % Diagonal conversions
    \draw[conv arrow, dashed, UofSC50Black] (de.south east) -- (bd.north west);
    \draw[conv arrow, dashed, UofSC50Black] (tf.south west) -- (ss.north east);

    % Central note
    \node[draw=UofSCGarnet, fill=UofSCSandstorm, font=\footnotesize,
          align=center, sharp corners, inner sep=4pt] at (0, 0) {
        All forms describe\\the same system
    };

\end{tikzpicture}
\caption{Conversion pathways between the four main system representations. Each representation offers different insights: differential equations show physical relationships, transfer functions simplify frequency analysis, state-space enables modern control design, and block diagrams provide visual system structure.}
\label{fig:representation_conversion}
\end{figure}
